\begin{figure}[htbp]
\centering
\begin{tcolorbox}[
    colback=gray!5!white,
    colframe=gray!60!black,
    title=\textbf{Prompt for AAC},
    fonttitle=\sffamily\bfseries,
    boxrule=0.8pt,
    arc=3pt,
    left=8pt, right=8pt, top=8pt, bottom=8pt
]
\small

You are provided with four images, consisting of two pairs of image editing data. In each pair, the first image is the original, and the second is the result generated by an image generation model. This model was trained on such pairs using a generic prompt. Your task is to refine and enrich this prompt to more accurately describe the transformation applied by this effect while maintaining its generality.

\vspace{0.5em}
\noindent\textbf{Constraints and Requirements:}
\begin{enumerate}
    \item \textbf{Subject Generalization}: Do not use specific descriptions for the subjects. Instead, use generic categories (e.g., human, pet, animal) based on the provided images.
    \item \textbf{Perspective Consistency}: Observe the viewpoint/perspective of the editing results. Describe the perspective shift only if it is consistently applied across all pairs; otherwise, omit it.
    \item \textbf{Reference Baseline}: Refer to the provided rough prompt: \texttt{"\{\}"}.
    \item \textbf{Output Specification}: Directly return the refined prompt without any supplementary or conversational text.
\end{enumerate}

\end{tcolorbox}
\caption{The prompt template used for refining and enriching the generic editing prompt based on visual samples and a baseline description.}
\label{fig:aac_prompt}
\end{figure}